% 14-vra.tex — Chapter 14: Value Range Analysis
\chapter{Value Range Analysis}
\label{chap:14-vra}
\index{VRA}
\index{value range analysis}

\pnum
Value Range Analysis (VRA) tracks compile-time integer value ranges for local
variables throughout a function body.  It is the core mechanism by which the
compiler statically proves array accesses safe without runtime checks.  VRA is
implemented in \srcref{src/sema/ranges.h}.

% -------------------------------------------------------------------------
\section{Range representation}
\label{sec:vra-range}
\index{Range struct}
% -------------------------------------------------------------------------

\pnum
A range is a closed integer interval $[\min, \max]$:

\begin{ccode}
typedef struct {
    int64_t min;
    int64_t max;
    bool    known;
} Range;
\end{ccode}

\pnum
\cname{known = false} means the range is unknown (unbounded).  Arithmetic on an
unknown range propagates the unknown: any operation whose operand is unknown
returns \cname{range\_unknown()}.  Three constructors:

\begin{center}
\begin{tabular}{ll}
\toprule
\textbf{Constructor} & \textbf{Produces} \\
\midrule
\cname{range\_unknown()} & Unknown (known=false, min=max=0) \\
\cname{range\_const(v)}  & $[v, v]$ — a known singleton \\
\cname{range\_make(lo, hi)} & $[\textit{lo}, \textit{hi}]$ \\
\bottomrule
\end{tabular}
\end{center}

% -------------------------------------------------------------------------
\section{Saturating arithmetic}
\label{sec:vra-sat}
\index{saturating arithmetic}
% -------------------------------------------------------------------------

\pnum
All range arithmetic uses saturating 64-bit operations to prevent silent
integer overflow.  Four helpers:

\begin{center}
\begin{tabular}{ll}
\toprule
\textbf{Helper} & \textbf{Saturates at} \\
\midrule
\cname{sat\_add\_i64(a,b)} & \texttt{INT64\_MAX}/\texttt{INT64\_MIN} \\
\cname{sat\_sub\_i64(a,b)} & \texttt{INT64\_MAX}/\texttt{INT64\_MIN} \\
\cname{sat\_mul\_i64(a,b)} & \texttt{INT64\_MAX}/\texttt{INT64\_MIN}, handles \texttt{INT64\_MIN/-1} UB \\
\cname{sat\_div\_i64(a,b)} & guards \texttt{INT64\_MIN/-1} UB, returns 0 for \texttt{b=0} \\
\bottomrule
\end{tabular}
\end{center}

% -------------------------------------------------------------------------
\section{RangeTable}
\label{sec:vra-table}
\index{RangeTable}
% -------------------------------------------------------------------------

\pnum
A \cname{RangeTable} maps variable \cname{Id *} to \cname{Range}.  It also
maintains a list of \cname{ConstraintEntry} records for linear constraints of
the form $x = y + c$:

\begin{itemize}
\item \cname{range\_set(table, id, range)} inserts or updates the range for a
  variable.
\item \cname{range\_get(table, id)} looks up a variable's range; returns
  \cname{range\_unknown()} if not found.
\item \cname{sema\_eval\_range(expr, table)} recursively evaluates an expression
  to a \cname{Range} by combining sub-expression ranges with saturating arithmetic.
\end{itemize}

% -------------------------------------------------------------------------
\section{Constraint propagation}
\label{sec:vra-constraints}
\index{constraint propagation}
% -------------------------------------------------------------------------

\pnum
The \cname{STMT\_ASSIGN} handler in \cname{walk\_stmt} adds linear constraints
to the table when the right-hand side has the form \lain{y + c}, \lain{c + y},
or \lain{y - c}:

\begin{ccode}
// x = y + c
constraint_add(table, x_id, y_id,  c);
constraint_add(table, y_id, x_id, -c);
\end{ccode}

\pnum
These bidirectional constraints allow the VRA to tighten bounds on either
variable when the other is constrained.  For example, if \texttt{y} is known
in range $[0, 7]$ and \texttt{x = y + 1}, the VRA can infer
\texttt{x} $\in [1, 8]$.

% -------------------------------------------------------------------------
\section{Condition-driven range refinement}
\label{sec:vra-cond}
\index{range refinement}
% -------------------------------------------------------------------------

\pnum
\cname{sema\_apply\_constraint(cond, table)} interprets a boolean condition as a
range restriction.  For example, \lain{if x < 10} causes the VRA to tighten
the upper bound of \texttt{x} to 9 in the then-branch.
\cname{sema\_apply\_negated\_constraint} applies the negation in the else-branch.

\pnum
Loop widening is applied before and after every loop body.
\cname{sema\_widen\_loop(body, table)} sets the range of any variable assigned
inside the loop to \cname{range\_unknown()}.  This is a conservative
over-approximation of the fixpoint; the VRA does not compute widening operators.

% -------------------------------------------------------------------------
\section{In-guard table}
\label{sec:vra-inguards}
\index{in-guard table}
\index{InGuardEntry}
% -------------------------------------------------------------------------

\pnum
The \lain{in} operator (e.g.\ \lain{idx in arr}) serves as a bounds-proving
condition.  When used as a loop condition or if-condition, it implies
$0 \leq \mathtt{idx} < \mathtt{arr.len}$.  The in-guard table records these
proven guards so the bounds checker can skip runtime checks for subscripts inside
the guarded region.

\begin{structdoc}{InGuardEntry}{src/sema.h:43}
\cname{index}     & \cname{Expr *}      & The index expression proven in-bounds. \\
\cname{container} & \cname{Expr *}      & The array/slice expression. \\
\cname{next}      & \cname{InGuardEntry *} & Next entry in the linked list. \\
\end{structdoc}

\pnum
\cname{sema\_push\_in\_guards(cond)} extracts every \lain{in} sub-expression from
a (possibly conjunctive) condition and adds an entry.  Guards are pushed at the
start of an \lain{if} or \lain{while} body and popped on scope exit.
\cname{sema\_is\_in\_guarded(index, container)} performs a structural-equality
search through the guard list.
